\scanpage{50}

2) Soit $k$ un corps, pas néc. alg. clos. Alors en général, un grp formel \inserted{$G$} sur $k$, même \inserted{une forme de} $G_{1,1}$ \unsure{pour $k$ fini}, ne peut se réaliser comme ss-grp formel d'un groupe formel associé à une VA définie sur $k$ --- et ceci même si on se permet une extension finie \deleted{de $k$} (ou de t.f.) de $k$. Prenons p.ex. $k$ corps fini, $E$ courbe elliptique hypersingulière sur $k$, \inserted{à multipl. compl. définie sur $k''$} \deleted{$\unsure{G}=\widehat E$} $G_0=\widehat E$, alors $\pi=\pi_1(k)$ opère trivialt sur $\operatorname{End}(E)$ et $\operatorname{End}(G_0)\deleted{\unsure{\mathcal A}}$ \inserted{$\simeq\operatorname{End}(E)\otimes_{\mathbb Z}\mathbb Z_p=\mathcal A$}, qui sont des \deleted{algèbres} \inserted{chacun un} ordres maximaux des algèbres de quaternions sur $\mathbb Z$ resp. $\mathbb Z_p$. On peut tordre donc $G_0$ par n'importe quelle unité $\theta$ de $\mathcal A$, \inserted{le nouveau frobénius devenant alors $\varphi'=\unsure{\Theta|\theta}\varphi$, $\varphi$ l'ancien frob.} \deleted{jouant le rôle de frobénius.} Mais si on veut que $G_0$ se plonge dans un $\widehat A$, $A$ une V.A., il faut par Weil que le polynôme caractéristique réduit de $\varphi'$ \inserted{$=\unsure{\Theta|\theta}\varphi$} \deleted{soit à} \inserted{ait comme racines des entiers alg.} \deleted{\unsure{dont et d'ailleurs tous les} coeff. ds $\mathbb Z$ --- (car ses val. propres doivent être des entiers algébriques),} \inserted{conjugués sont de val. absolue $q^{1/2}$, p.ex. il faut} \deleted{\unsure{et que} \illegible{2} \unsure{particulier il faut valeurs absolues vont être} $q^{1/2}$, \unsure{donc il faut} $\det\unsure{\varphi'}\quad\det\varphi'=1\det\varphi$,} que $\det\varphi'=\det\varphi\det\unsure{\Theta|\theta}$ \deleted{soit} soit un entier alg. dont \deleted{\unsure{toutes}} les conjugués sont de val. absolue \unsure{,} ce qui implique que $\det\unsure{\Theta|\theta}$ doit être une unité alg. dont tous les conjugués sont de val abs. $1$.

\margin{60}
